\documentclass[11pt]{article}
\usepackage[T1]{fontenc}
\usepackage{lmodern}
\usepackage[margin=1.05in]{geometry}
\usepackage{amsmath,amssymb,amsthm,mathtools,mathrsfs}
\usepackage{microtype,booktabs,longtable,array}
\usepackage{xurl}
\usepackage[colorlinks=true,linkcolor=blue,citecolor=blue,urlcolor=blue]{hyperref}
\setlength{\emergencystretch}{3em}
\newtheorem{theorem}{Theorem}[section]
\newtheorem{proposition}[theorem]{Proposition}
\newtheorem{lemma}[theorem]{Lemma}
\theoremstyle{definition}
\newtheorem{definition}[theorem]{Definition}
\theoremstyle{remark}
\newtheorem{remark}[theorem]{Remark}
\theoremstyle{plain}
\newtheorem{mainthm}{Theorem}
\renewcommand{\themainthm}{\Alph{mainthm}}
\newcommand{\Q}{\mathbb Q}
\newcommand{\Z}{\mathbb Z}
\newcommand{\N}{\mathbb N}
\newcommand{\R}{\mathbb R}
\newcommand{\C}{\mathbb C}
\newcommand{\OO}{\mathcal O}
\newcommand{\wt}[1]{\lvert #1\rvert_w}
\newcommand{\Emb}{\operatorname{Emb}}
\DeclareMathOperator{\lcm}{lcm}
\hypersetup{pdftitle={The irrationality exponents of log(2) + sqrt(2) and sqrt(2) log(2) are at most 4},pdfauthor={Ryan Matthew Casper},pdfsubject={Affine logarithms, rational approximation, and a degree-four norm argument}}
\title{The irrationality exponents of\newline
\(\log(2)+\sqrt2\) and \(\sqrt2\log(2)\)\newline are at most \(4\)}
\author{Ryan Matthew Casper}
\date{11 October 2026\\\small Preprint working draft}
\begin{document}
\maketitle

\begin{abstract}
We prove that \(\log 2+\sqrt2\), \(\sqrt2\log 2\), and
\(\log(3+2\sqrt2)/\sqrt2+\sqrt3\) are irrational and have rational
irrationality exponent at most four. For each of these numbers \(x\) and
each real \(\nu>4\), the proof gives an eventual lower bound
\(q^{-\nu}\le |x-p/q|\), uniformly in the integer numerator and positive
integer denominator; the fractions need not be reduced. The common argument
takes place in a number field of degree four. Two embeddings satisfy an
affine exponential identity and receive a logarithmic collision estimate;
the other two receive a coefficient estimate. All four evaluate the same
cleared algebraic-integer determinant. Its norm supplies the arithmetic
lower bound. We state the inherited weighted interpolation and collision
inputs explicitly, explain the parameter choices and denominator clearing,
and give both branches of the determinant contradiction. The accompanying
Lean package includes the three concrete bounds, their irrationality and
exponent consequences, and checked comparisons with the stated expressions.
The conclusions are upper bounds; they do not determine the exact exponents.
\end{abstract}

\section{Introduction and main results}\label{sec:intro}
Adding an algebraic number to a logarithm, or multiplying the logarithm by
one, changes the rational approximation problem. An approximation to
\(\log 2+\sqrt2\) becomes an approximation to \(\log 2\) by an element
of \(\Q(\sqrt2)\), rather than by a rational number. The same issue
appears for \(\sqrt2\log2\). We keep that algebraic arithmetic inside
the determinant construction and prove the following three statements.

Throughout the paper, \(\log\) denotes the natural real logarithm on
positive arguments, and square roots denote positive real square roots.
The symbol \(\mu\) denotes the rational irrationality exponent, with the
convention given in Section~\ref{sec:conventions}.

\begin{mainthm}[An algebraic shift]\label{thm:shift}
The number \(x_1=\log2+\sqrt2\) is irrational and
\[
 \mu(\log2+\sqrt2)\le4.
\]
\end{mainthm}

\begin{mainthm}[An algebraic multiple]\label{thm:scale}
The number \(x_2=\sqrt2\log2\) is irrational and
\[
 \mu(\sqrt2\log2)\le4.
\]
\end{mainthm}

\begin{mainthm}[A shifted quadratic period]\label{thm:period}
The number
\[
 x_3=\frac{\log(3+2\sqrt2)}{\sqrt2}+\sqrt3
\]
is irrational and \(\mu(x_3)\le4\).
\end{mainthm}

For each of the three numbers, we prove the more precise statement
\begin{equation}\label{eq:eventual}
 \forall\nu>4\quad\exists Q\in\N,\ Q\ge2:\quad
 \forall p\in\Z\ \forall q\in\N,\quad
 q\ge Q\ \Longrightarrow\ q^{-\nu}\le\left|x-\frac pq\right|.
\end{equation}
The threshold depends on \(x\) and \(\nu\), but not on the numerator
or denominator. There is no coprimality condition in
\eqref{eq:eventual}. This detail also lets the bound prove irrationality:
an exact rational representation could otherwise be repeated at arbitrarily
large denominators.

The three applications use a common theorem in a degree-four number field.
At two embeddings \(\sigma\), integral elements \(a,u,v\) satisfy
\[
 \exp\bigl(\sigma(u)(x-\sigma(v))\bigr)=\sigma(a).
\]
These are the embeddings where good rational approximations to \(x\)
produce small translation errors. At the other embeddings we need only a
bound on the matrix coefficients. Interpolation selects a nonzero minor at
one positive real embedding. We clear that minor over the ring of integers,
keep its column selection fixed, and take its norm.

The factor four in the final exponent comes from this comparison. There
are four copies of the arithmetic clearing cost and two copies of the full
analytic saving. After division by two, the leading arithmetic cost is
\(2(1-b_H)\), where \(b_H\) is the mean transverse row weight. The
parameters are chosen to defeat this cost at every exponent \(\nu>4\).
The coarser estimates still have to be paid for; increasing the determinant
height does not remove all of them. We therefore choose the dimension and
the approximation denominators before letting that height grow.

Our starting point is the weighted logarithmic interpolation and collision
framework in OpenAI's \texttt{PiExponent} library~\cite{openai}, together
with the rational-logarithm and quadratic-logarithm adaptations
\cite{rational,quadratic}. The present argument uses those mathematical
inputs, including their algebraic geometry and complex analysis. It does
not provide a separate proof of their foundations. Sections~\ref{sec:geometry}
and~\ref{sec:analytic} state the interfaces we need, so the dependence is
visible before either input is applied.

\subsection{A targeted comparison with earlier work}
Luchin and Salikhov~\cite{luchin} prove that, for every
\(s>7.98101\), sufficiently large coefficient height
\(P=\max_i|p_i|\) gives
\begin{equation}\label{eq:LS}
 \left|\log2-\frac{p_1\sqrt2+p_2}{p_3\sqrt2+p_4}\right|>P^{-s},
 \qquad (p_3,p_4)\ne(0,0),
\end{equation}
where the \(p_i\) are integers. This is Theorem~1 on page~549 of their
paper. Two direct substitutions make the relation to our first two
numbers precise.

For \(x_1\), take \((p_1,p_2,p_3,p_4)=(-q,p,0,q)\). The left side of
\eqref{eq:LS} is then \(|x_1-p/q|\). For \(x_2\), take
\((p_1,p_2,p_3,p_4)=(0,p,q,0)\); the left side is
\(|x_2-p/q|/\sqrt2\). In either case \(P=\max(|p|,q)\). If
\(|x_j-p/q|<1\), then, since these two targets are positive,
\[
 q\le P\le (|x_j|+1)q.
\]
Given \(\nu>7.98101\), choose \(7.98101<s<\nu\). The fixed
multiplicative constants in \eqref{eq:LS} are absorbed by
\(q^{\nu-s}\) for sufficiently large \(q\); errors at least one
already satisfy the desired lower bound. Thus that theorem implies
\(\mu(x_1),\mu(x_2)\le7.98101\). There is no squaring of the
denominator height in this substitution.

Theorems~\ref{thm:shift} and~\ref{thm:scale} improve this particular
consequence to four. This is a comparison with a specified theorem, not an
exhaustive survey or a claim of priority. We make no corresponding
best-known-bound claim for \(x_3\).

\section{Approximation conventions}\label{sec:conventions}
For an irrational real number \(x\), let
\[
 \mathcal A(x)=\left\{\nu>0:
 \begin{array}{l}
 \text{infinitely many distinct reduced fractions }p/q,\ q\ge2,\\[-2pt]
 \text{satisfy }0<|x-p/q|<q^{-\nu}
 \end{array}\right\}.
\]
We use \(\mu(x)=\sup\mathcal A(x)\), allowing the value \(+\infty\).
For the numbers considered here the set will be nonempty and bounded
above, so the ordinary real supremum suffices. Dirichlet approximation
gives \(2\in\mathcal A(x)\) for irrational \(x\).

\begin{lemma}\label{lem:exponent}
If a real number \(x\) satisfies \eqref{eq:eventual}, then \(x\) is
irrational and \(2\le\mu(x)\le4\).
\end{lemma}
\begin{proof}
Suppose \(x=r/s\), with \(r\in\Z\) and \(s\) a positive integer.
Apply \eqref{eq:eventual} with \(\nu=5\). For an integer \(N\)
large enough that \(Ns\ge Q\), put \(p=Nr\) and \(q=Ns\).
The approximation error is zero, contradicting \(q^{-5}>0\).

Now fix \(\nu>4\). Every reduced fraction counted in
\(\mathcal A(x)\) at exponent \(\nu\) has \(q<Q\). Its error is
less than one, hence \(|p|<q(|x|+1)\). There are only finitely many
such denominators and numerators. Therefore no \(\nu>4\) belongs to
\(\mathcal A(x)\), which proves the upper bound. Dirichlet's theorem
gives the lower bound.
\end{proof}

We will prove \eqref{eq:eventual} by excluding an unbounded supply of
denominators for the opposite estimate. More precisely, for each fixed
\(\nu>4\) we rule out
\begin{equation}\label{eq:bad}
 \forall Q\in\N\quad\exists p\in\Z\ \exists q\in\N,\quad
 q\ge Q,\qquad \left|x-\frac pq\right|\le q^{-\nu}.
\end{equation}
Its negation gives a threshold beyond which the error is strictly greater
than \(q^{-\nu}\). Increasing that threshold to at least two and
weakening the inequality gives \eqref{eq:eventual}. We use the non-strict
inequality in \eqref{eq:bad} because it is convenient for selecting
approximation data; the final exponent convention still uses a strict
inequality.

\section{The common degree-four theorem}\label{sec:common}
For a number field \(F\), write \(\OO_F\) for its ring of integers and
\(\Emb(F)=\operatorname{Hom}_{\Q\text{-alg}}(F,\C)\). Embeddings are
counted as maps. In particular, two distinct embeddings can have the same
value on \(a,u,v\).

\begin{theorem}[Two compatible embeddings]\label{thm:common}
Let \(F\) be a number field with \([F:\Q]=4\), and let
\(a,u,v\in\OO_F\). Suppose there is an embedding \(\sigma_0\)
such that \(\sigma_0(a)=b_0\) is real, \(b_0>0\), and \(b_0\ne1\).
Let \(x\in\R\setminus\{0\}\). Suppose a set
\(G\subset\Emb(F)\) of exactly two embeddings satisfies
\begin{equation}\label{eq:compatible}
 \exp\bigl(\sigma(u)(x-\sigma(v))\bigr)=\sigma(a),
 \qquad |\sigma(u)|\le2\qquad(\sigma\in G).
\end{equation}
Then \(x\) satisfies \eqref{eq:eventual}, and consequently is irrational
with \(\mu(x)\le4\).
\end{theorem}

The distinguished embedding \(\sigma_0\) is used for interpolation.
It need not belong to \(G\). The theorem imposes no exponential
compatibility at the two remaining embeddings, and they need not be real.
Complex absolute values are used throughout the common argument. We call
the members of \(G\) the \emph{good embeddings}, meaning only that
\eqref{eq:compatible} holds there.

\subsection{Fixed bounds for the embeddings}
Fix \(\nu>4\) and assume \eqref{eq:bad}. Before choosing any rational
approximations, put
\begin{align}
 \kappa&=\sum_{\sigma\in\Emb(F)}\log\max(1,|\sigma(a)|),\label{eq:budgets}\\
 R&=\sum_{\sigma\in\Emb(F)}|\sigma(u)|\bigl(|x|+1+|\sigma(v)|\bigr),
 &R'&=2R+2,\notag\\
 \rho&=100\max\left(1,|x|,
       \sum_{\sigma\in\Emb(F)}|\sigma(u)(x-\sigma(v))|\right).\notag
\end{align}
All summands are nonnegative. Thus \(\kappa,R\ge0\),
\(\rho\ge100\), \(100|x|\le\rho\), and, separately for every embedding,
\begin{equation}\label{eq:budgetbounds}
 \begin{split}
 \log\max(1,|\sigma(a)|)&\le\kappa,\\
 |\sigma(u)|(|x|+1+|\sigma(v)|)&\le R,\\
 100|\sigma(u)(x-\sigma(v))|&\le\rho.
 \end{split}
\end{equation}
These deliberately generous bounds keep the later choices uniform over
the embeddings. The enlargement from \(R\) to \(R'\) will pay for the
fixed multiplier in the affine approximation error.

\section{Choosing the parameters}\label{sec:parameters}
There are two quite different limiting operations in the proof. First we
choose a dimension and a finite set of good rational approximations.
Then we hold them fixed and increase the matrix height \(H\). Keeping
that order matters: a term such as \(1/w_*\), where \(w_*\) is the
smallest approximation weight, does not tend to zero with \(H\).

Set
\begin{equation}\label{eq:constants}
 \Lambda=\log4+4,\qquad c=\frac{\log2}{4}.
\end{equation}
The first is an admissible constant for the least-common-multiple estimate;
the second is the collision constant. Choose rational shape parameters
\begin{equation}\label{eq:shape}
 0<\theta<A<B<1,\quad C>1,\quad B<C^{-1},\quad B<C\theta<1,
\end{equation}
and a real \(0<\eta<1\), so that
\begin{equation}\label{eq:gap}
 g=\nu\bigl(A(1-\eta)-\theta\bigr)-2(1-\theta)>0.
\end{equation}
The inherited shape lemma supplies these parameters by applying its
exponent-two construction to \(\nu/2>2\). Only the shape construction
uses the rescaled exponent. Our rational approximations retain exponent
\(\nu\).

Choose \(0<\epsilon<g\) with \(\epsilon\le1/2\), and then choose
\(F_0>2/\theta\) such that \(\nu/F_0<\epsilon/3\). For a sufficiently
large integer \(m\ge1\), set
\begin{equation}\label{eq:dimension}
 K=\lfloor C^m\rfloor,\qquad w_0=B^{-m},\qquad
 v_0=2K\theta^m w_0.
\end{equation}
Here \(K\ge1\), the weights \(w_0,v_0\) are positive rationals, and
\begin{equation}\label{eq:volume}
 K\frac{w_0}{v_0}\theta^m=\frac12,
 \qquad K\theta^m<1.
\end{equation}
We choose \(m\) to make
\begin{equation}\label{eq:dimensionmargin}
 \frac{2\Lambda F_0m+3\log2}{v_0}
 +\frac{(\rho+\kappa)K+\log(3/2)}{w_0}<\frac\epsilon3
\end{equation}
and simultaneously
\begin{equation}\label{eq:collisionlimit}
 C_\infty:=\frac{c\eta^2K\theta^m}{(m+1)v_0A^m}>3.
\end{equation}

Why can these requirements be met together? With the shape fixed,
\(K/w_0\le(CB)^m\to0\), while \(v_0\) grows exponentially because
\(C\theta/B>1\). For example, once \(C^m\ge2\),
\(K\ge C^m/2\), so \(v_0\ge(C\theta/B)^m\). Thus even
\(m/v_0\to0\). On the other hand, substitution from
\eqref{eq:dimension} gives the exact identity
\[
 C_\infty=\frac{c\eta^2}{2(m+1)}\left(\frac BA\right)^m,
\]
which tends to infinity. The floor in \(K\) causes no difficulty in
either comparison.

Choose a small positive rational \(\zeta\) such that
\begin{equation}\label{eq:slack}
 \frac{(1+3\zeta)^{m+1}}2<1,\qquad
 (1+3\zeta)^mK\theta^m<1,\qquad (1+\zeta)\theta<1.
\end{equation}
We use \(\zeta\) for this slack to reserve \(\sigma\) for field
embeddings. All three inequalities follow by taking \(\zeta\) close
enough to zero, since their zero-slack versions are strict.

\subsection{Successive approximation weights}
The hypothesis \eqref{eq:bad} supplies a sequence of approximations whose
denominators can be required to exceed any previously specified threshold.
Choose its first \(m\) members \(p_i/q_i\), \(1\le i\le m\), with
\begin{equation}\label{eq:approxdata}
 q_i\ge2,\quad p_i\ne0,\quad
 |x-p_i/q_i|\le q_i^{-\nu},\qquad
 w_i=\lceil\log q_i\rceil\ge1.
\end{equation}
The condition \(p_i\ne0\) is available because \(x\ne0\): zero
numerators cannot satisfy the approximation inequality once \(q_i\)
is sufficiently large. Let \(w_* =\min_i w_i>0\).

The weights must also be sufficiently separated for interpolation. Here
is the precise finite condition we impose. Set
\[
 d=(w_0,w_1,\ldots,w_m),\qquad
 e=(v_0,w_1/\theta,\ldots,w_m/\theta),
\]
and
\begin{equation}\label{eq:separationconstant}
 D_m(\zeta)=2(m+2)^{m+2}
       \left(1+\frac{m+2}{\zeta}\right)^{m+2}.
\end{equation}
For equal-cardinality subsets \(P,Q\subseteq\{0,\ldots,m\}\), whenever
\(i>0\) belongs to \(P\setminus Q\) and membership in \(P,Q\)
agrees at every index above \(i\), require
\begin{equation}\label{eq:separation}
 D_m(\zeta)\prod_{j\in Q}e_j<\prod_{j\in P}d_j.
\end{equation}
After common indices are cancelled, the largest remaining index is on
the right. Its weight can dominate the finitely many products of earlier
weights. Successively increasing the next denominator therefore achieves
these inequalities. Distinct denominators alone would not be enough.

Finally choose the initial denominator threshold large enough to make the
remaining inverse-weight errors small. Define
\begin{align}
 E&=\frac{\Lambda F_0m}{v_0}+\Lambda\sum_{i=1}^m\frac1{w_i}
                          +\frac\theta{w_*},\label{eq:errors}\\
 U(S)&=\frac{\log2}{v_0}+\frac{\kappa K+\log(3/2)}{w_0}
                   +\frac{\log(2(KS+2))}{w_*}\quad(S\ge0),\notag\\
 T&=\frac\nu{F_0}+\frac{\log2}{v_0}
                    +\frac{\log4+\log(2K)+\nu}{w_*},\notag\\
 J&=\frac{\rho K}{w_0}+\frac{\log2}{v_0}
                    +\frac{\log(200K)+\log(\rho/100)}{w_*},
 &V&=T+J.\notag
\end{align}
The part of \(2E+U(R')+V\) independent of \(w_i\) is
\(\nu/F_0\) plus the left side of \eqref{eq:dimensionmargin}.
The rest is
\begin{equation}\label{eq:weightmargin}
 2\Lambda\sum_i\frac1{w_i}
 +\frac{2\theta+\log(2(KR'+2))+\log4+\log(2K)+\nu
       +\log(200K)+\log(\rho/100)}{w_*}.
\end{equation}
Since \(\sum_i1/w_i\le m/w_*\), an arbitrarily large initial weight
makes \eqref{eq:weightmargin} less than \(\epsilon/3\), without
interfering with later separation. We have proved the two margins
\begin{equation}\label{eq:twomargins}
 2E+U(R')+V<g,
 \qquad 2+2E+U(R')+V<C_\infty.
\end{equation}
For the second inequality, the error sum is less than
\(\epsilon\le1/2\), whereas \(C_\infty>3\).

The order of choice is now complete: shape, \(F_0\), dimension, slack,
initial approximation weight, successive approximations. From this point
on, every quantity in \eqref{eq:errors} is fixed. Only \(H\), the
matrix indices, and the selected minor will vary.

\section{Weighted interpolation and the matrix}\label{sec:geometry}
We first state the inherited interpolation input in a form that separates
its numerical hypotheses from its conclusion. A polynomial in
\(Y,X_1,\ldots,X_m\) has weighted degree at most \(H\), with degree
weights \(d_i\), if every monomial \(Y^hX^\gamma\) in its support
satisfies \(d_0h+\sum_i d_i\gamma_i\le H\). A logarithmic jet packet
with weights \(e_i\) consists of the coefficients of
\(t^sZ^\beta\) with \(e_0s+\sum_i e_i\beta_i<H\), after the formal
substitution
\begin{equation}\label{eq:formalcoords}
 Y=Y_j(1+t),\qquad X_i=c_{ji}+\log(1+t)+Z_i.
\end{equation}
Here the logarithm is the formal power series at zero.

\begin{proposition}[Inherited weighted interpolation]\label{prop:interpolation}
Let \(d_i,e_i\) be positive rational weights, \(\zeta>0\) rational,
\(Y_0,\ldots,Y_{K-1}\) distinct nonzero complex numbers, and
\(c_{ji}\) arbitrary complex numbers. Suppose
\begin{equation}\label{eq:geomvolume}
 K(1+3\zeta)^{m+1}\frac{\prod_{i=0}^m d_i}{\prod_{i=0}^m e_i}<1,
 \qquad (1+\zeta)d_i<e_i\quad(1\le i\le m),
\end{equation}
and the separated-product condition \eqref{eq:separation} holds for these
weights. There is a positive rational scale \(\mathcal R\) such that,
for every sufficiently large integer \(n\), the map from polynomials of
weighted degree at most \(n\mathcal R\) to their logarithmic jet packets
of weight strictly below \(n\mathcal R\) at all \(K\) centers is
surjective.
\end{proposition}

This proposition is the weighted logarithmic interpolation theorem of the
inherited source~\cite{rational}. Its proof passes through ampleness of
an interpolation bundle and eventual restriction to jets, using the
weighted geometric results in~\cite{openai}. We use that theorem with
its stated hypotheses. In particular, the arbitrary additive centers in
\eqref{eq:formalcoords} do not carry a hidden Diophantine assumption.

\subsection{The finite matrix}
For \(H>0\), define the column and row sets
\begin{align}
 \mathcal C_H&=\{(h,\gamma)\in\N\times\N^m:
                   w_0h+\wt\gamma\le H\},\label{eq:indices}\\
 \mathcal R_H&=\{(j,s,\beta):0\le j<K, s\in\N,\ \beta\in\N^m,
                   v_0s+\wt\beta/\theta<H\},\notag
\end{align}
where \(\wt\alpha=\sum_iw_i\alpha_i\), and put \(M=M(H)=\#\mathcal R_H\).
This is a positive integer because \((0,0,0)\) is a row. Set
\begin{equation}\label{eq:truncations}
 T_i=\left\lceil\frac{F_0w_i}{v_0}\right\rceil,\qquad
 L_i(t)=\sum_{1\le k<T_i}\frac{(-1)^{k+1}}k t^k.
\end{equation}
For an embedding \(\sigma\), write
\begin{equation}\label{eq:centers}
 z_\sigma=\sigma(a),\qquad
 r_{\sigma i}=\sigma(u)\left(\frac{p_i}{q_i}-\sigma(v)\right).
\end{equation}
Our matrix, with rows and columns from \eqref{eq:indices}, is
\begin{equation}\label{eq:matrix}
 \mathcal M_\sigma(H)_{(j,s,\beta),(h,\gamma)}
 =z_\sigma^{jh}[t^sZ^\beta]\,
       (1+t)^h\prod_{i=1}^m
       \bigl(Z_i+jr_{\sigma i}+L_i(t)\bigr)^{\gamma_i}.
\end{equation}
The factor \(z_\sigma^{jh}\) is part of the matrix. It contributes to
both the arithmetic and analytic estimates and cannot be omitted.

\begin{lemma}\label{lem:cofinal}
The matrix \(\mathcal M_{\sigma_0}(H)\) has full row rank for a cofinal
set of heights \(H>0\).
\end{lemma}
\begin{proof}
Apply Proposition~\ref{prop:interpolation} with the weights \(d,e\)
chosen above, multiplicative centers \(Y_j=b_0^j\), and additive
centers \(c_{ji}=jr_{\sigma_0 i}\). Since \(b_0>0\) and
\(b_0\ne1\), these multiplicative centers are distinct and nonzero.
The volume ratio is \(1/2\) by \eqref{eq:volume}, the coordinate
inequalities follow from \((1+\zeta)\theta<1\), and separation is
\eqref{eq:separation}.

To pass to \eqref{eq:matrix}, let
\(\tau_i(t)=\log(1+t)-L_i(t)\). It has order at least \(T_i\), and
\[
 v_0T_i\ge F_0w_i> w_i/\theta.
\]
Thus \(Z_i\mapsto Z_i+\tau_i(t)\) and its inverse preserve the
weighted filtration with weights \(v_0,w_i/\theta\). They induce
inverse linear maps on the finite jet packet spaces. Applying this
substitution to the truncated expression recovers the full logarithmic
expression. Consequently the full and truncated packet maps are
surjective at exactly the same heights. Individual coefficients can
change; it is the invertible change of packet coordinates that preserves
surjectivity.

In the monomial bases, the truncated packet map has precisely the entries
\eqref{eq:matrix}. Proposition~\ref{prop:interpolation} therefore gives
full row rank for all sufficiently large heights of the form
\(H=n\mathcal R\). Since \(\mathcal R>0\), these heights are cofinal.
\end{proof}

At each such height, choose a nonzero full-row minor and denote its column
selection by \(s_H\). Use this \emph{same} column selection at every
embedding and let
\begin{equation}\label{eq:minor}
 \Delta_\sigma=\det\bigl(\mathcal M_\sigma(H)|_{s_H}\bigr).
\end{equation}
At this stage only \(\Delta_{\sigma_0}\ne0\) is known. The common
integral lift below will show that all four determinants are nonzero.

\section{Clearing the selected minor}\label{sec:arithmetic}
Let \(e_i^H=\lfloor H/w_i\rfloor\) and
\(D_i=\lcm(1,\ldots,T_i)\). Every selected column satisfies
\(\gamma_i\le e_i^H\). Multiply row \((j,s,\beta)\) by
\(\prod_iD_i^{e_i^H}q_i^{-\beta_i}\), and selected column
\((h,\gamma)\) by \(\prod_iq_i^{\gamma_i}\). Their total determinant
factor is the positive rational number
\begin{equation}\label{eq:clearing}
 \mathscr C_H=
 \left(\prod_iD_i^{e_i^H}\right)^M
 \left(\prod_{r\in\mathcal R_H}\prod_iq_i^{-\beta_{ri}}\right)
 \left(\prod_{c\in s_H}\prod_iq_i^{\gamma_{ci}}\right).
\end{equation}
This factor need not be an integer. The assertion we need is that the
scaled entries, and therefore their determinant, are algebraic integers.

\begin{lemma}[A common integral lift]\label{lem:integral}
There is an element \(Z_H\in\OO_F\setminus\{0\}\) such that
\begin{equation}\label{eq:lift}
 \sigma(Z_H)=\mathscr C_H\Delta_\sigma
 \qquad\text{for every }\sigma\in\Emb(F).
\end{equation}
In particular, \(\Delta_\sigma\ne0\) at every embedding.
\end{lemma}
\begin{proof}
Define
\[
 n_i=u(p_i-q_iv)\in\OO_F.
\]
Then \(r_{\sigma i}=\sigma(n_i)/q_i\). If some
\(\beta_i>\gamma_i\), the entry in \eqref{eq:matrix} is zero. In
the remaining case, extracting \(Z^\beta\) gives the factor
\(\prod_i\binom{\gamma_i}{\beta_i}\) and powers with exponents
\(d_i'=\gamma_i-\beta_i\ge0\). After the row and column factors
have been applied, the entry is the image under \(\sigma\) of
\begin{equation}\label{eq:entrylift}
 a^{jh}\left(\prod_i\binom{\gamma_i}{\beta_i}\right)
 [t^s]\,(1+t)^h
 \prod_i D_i^{e_i^H}\bigl(jn_i+q_iL_i(t)\bigr)^{d_i'}.
\end{equation}
Each polynomial \(D_i(jn_i+q_iL_i(t))\) has coefficients in
\(\OO_F\), because \(D_i\) clears every denominator occurring in
\(L_i\). Also \(d_i'\le\gamma_i\le e_i^H\), so
\[
 D_i^{e_i^H}(jn_i+q_iL_i)^{d_i'}
 =D_i^{e_i^H-d_i'}\bigl(D_i(jn_i+q_iL_i)\bigr)^{d_i'}
\]
is integral coefficientwise. The other factors in \eqref{eq:entrylift}
are integral as well.

Use these entries to form one square matrix over \(\OO_F\), and let
\(Z_H\) be its determinant. Determinants commute with embeddings,
which proves \eqref{eq:lift}. Since \(\mathscr C_H>0\) and
\(\Delta_{\sigma_0}\ne0\), we have \(Z_H\ne0\). Every embedding
of a field is injective, so \(\sigma(Z_H)\ne0\) for every \(\sigma\).
Equation~\eqref{eq:lift} then gives the remaining assertion.
\end{proof}

Define the normalized mean transverse row weight by
\begin{equation}\label{eq:mean}
 b_H=\frac1{MH}\sum_{r\in\mathcal R_H}\wt{\beta_r}.
\end{equation}
The row condition gives \(0\le b_H\le\theta\).

\begin{lemma}[Clearing cost]\label{lem:clearingcost}
For every \(H>0\) and selected full-row minor,
\begin{equation}\label{eq:clearingcost}
 \log\mathscr C_H\le MH(1-b_H+E).
\end{equation}
\end{lemma}
\begin{proof}
The inherited elementary least-common-multiple estimate is
\(\log D_i\le\Lambda T_i\). Since
\(T_i\le F_0w_i/v_0+1\) and \(e_i^H\le H/w_i\),
\[
 \sum_i e_i^H\log D_i
 \le H\left(\frac{\Lambda F_0m}{v_0}
                         +\Lambda\sum_i\frac1{w_i}\right).
\]
For the column factors, \(\log q_i\le w_i\) and the column degree
bound imply
\(\sum_i\gamma_{ci}\log q_i\le H\) for each selected column.
For the inverse row factors, \(\log q_i\ge w_i-1\) gives
\[
 -\sum_{r,i}\beta_{ri}\log q_i
 \le -MHb_H+\sum_{r,i}\beta_{ri}
 \le -MHb_H+\frac{MH\theta}{w_*}.
\]
Adding these three estimates in the logarithm of \eqref{eq:clearing}
proves the claim.
\end{proof}

The inverse row factors are useful arithmetic, not just a normalization.
They give the term \(-b_H\). Dropping them would lose the comparison
with the transverse weight in the analytic estimate.

\subsection{The norm inequality}
The nonzero algebraic integer \(Z_H\) has nonzero integer norm. Thus
\[
 1\le |N_{F/\Q}(Z_H)|=\prod_{\sigma\in\Emb(F)}|\sigma(Z_H)|,
\]
and \eqref{eq:lift} yields
\begin{equation}\label{eq:norm}
 0\le4\log\mathscr C_H+
              \sum_{\sigma\in\Emb(F)}\log|\Delta_\sigma|.
\end{equation}
All logarithms are legitimate by Lemma~\ref{lem:integral}.

\begin{lemma}[Normalized comparison]\label{lem:normcomparison}
Suppose \(\log|\Delta_\sigma|/(MH)\le B_H\) at the two good
embeddings and \(\log|\Delta_\sigma|/(MH)\le D_H\) at the two
remaining embeddings. Then
\begin{equation}\label{eq:normnormalized}
 -2(1-b_H)-2E-D_H\le B_H.
\end{equation}
\end{lemma}
\begin{proof}
Insert Lemma~\ref{lem:clearingcost} and the four determinant estimates
into \eqref{eq:norm}. The resulting upper bound on its right side is
\(4MH(1-b_H+E)+2MHB_H+2MHD_H\), and is nonnegative. Divide by
\(2MH>0\) and rearrange.
\end{proof}

This is where the degree of the field and the number of good embeddings
enter together. We have not bounded each \(|\sigma(Z_H)|\) below by
one. Such a bound would generally be false. The product is the integer
quantity.

\section{The analytic estimates}\label{sec:analytic}
Fix a good embedding \(\sigma\) and write
\[
 \omega=\sigma(u)(x-\sigma(v)),\qquad z=z_\sigma.
\]
Then \(e^\omega=z\) and \(100|\omega|\le\rho\). The translated
centers satisfy
\begin{align}
 r_{\sigma i}-\omega&=\sigma(u)(p_i/q_i-x),\label{eq:affineerror}\\
 |j(r_{\sigma i}-\omega)|
 &\le\exp\bigl(\log(2K)+\nu+\log2-\nu w_i\bigr),
 \qquad 0\le j<K.\notag
\end{align}
Indeed, \(|\sigma(u)|\le2\), the approximation error is at most
\(q_i^{-\nu}\), and \(\log q_i\ge w_i-1\). We have retained a
harmless factor two from the inherited approximation estimate, as well
as the extra \(\log2\) due to \(u\). We will explicitly absorb the
latter after combining the embeddings.

\subsection{The collision input}
For nonnegative integers \(h\), multi-indices \(\gamma,\alpha\), and
an integer \(j\), consider the formal coefficient
\[
 e^{j\omega h}[t^\ell Z^\alpha]
 (1+t)^h\prod_i\bigl(Z_i+j\omega+\log(1+t)\bigr)^{\gamma_i}.
\]
These are the entries to which the inherited collision estimate applies.
They are related to the entire functions
\begin{equation}\label{eq:columnfunction}
 f_{h,\gamma,\alpha}(z)
 =e^{hz}\left(\prod_i\binom{\gamma_i}{\alpha_i}\right)
           z^{\sum_i(\gamma_i-\alpha_i)}
 \quad(\alpha\le\gamma),
\end{equation}
and are zero when \(\alpha\not\le\gamma\). The logarithmic
coordinate is \(z=j\omega+\log(1+t)\). Repeated transverse
indices cause the corresponding rows to use the same family of entire
functions, which is the source of the collision saving.

\begin{proposition}[Inherited scaled collision estimate]\label{prop:collision}
Let a square matrix have \(M\ge1\) rows indexed by \(r\) and columns
indexed by \(c\). Assign a transverse multi-index \(\alpha_r\), an
integer \(0\le j_r<K\), and an order \(\ell_r\ge0\) to each row;
assign \((h_c,\gamma_c)\) with
\(w_0h_c+\wt{\gamma_c}\le H\) to each column. Suppose
\[
 H,w_0,v_0,w_*>0,\quad w_i\ge w_*,\quad
 \rho\ge100,\quad 100|\omega|\le\rho,
 \quad \sum_r\ell_r\le MH/v_0.
\]
Let \(n_\alpha=\#\{r:\alpha_r=\alpha\}\), and let \(B\) be the
matrix with entries
\begin{equation}\label{eq:collisionmatrix}
 B_{rc}=e^{j_r\omega h_c}[t^{\ell_r}Z^{\alpha_r}]
 (1+t)^{h_c}\prod_i
       \bigl(Z_i+j_r\omega+\log(1+t)\bigr)^{\gamma_{ci}}.
\end{equation}
Then
\begin{equation}\label{eq:collisionbound}
 |\det B|\le
 \exp\left(-c\sum_\alpha n_\alpha^2+
                  MH\bigl(J+r_{\rm coll}(M,H)\bigr)\right),
\end{equation}
where \(J\) is given in \eqref{eq:errors} and
\begin{equation}\label{eq:collisionremainder}
 r_{\rm coll}(M,H)
  =\frac{\log M+c-\log(1-2^{-1/2})}{H}.
\end{equation}
\end{proposition}

This is the scaled form of the inherited holomorphic collision estimate
\cite{openai,rational}. Its analytic domain is a disk of radius \(\rho K\).
The centers lie well inside this disk:
\(\lvert j\omega\rvert+3/4\le\rho K/2\). On its boundary, the
weighted degree of the columns bounds the entire functions in
\eqref{eq:columnfunction} by
\[
 \exp\left(H\left(\frac{\rho K}{w_0}
                           +\frac{\log(2\rho K)}{w_*}\right)\right).
\]
The derivative-order budget supplies \(\log2/v_0\), and the identity
\(\log(2\rho K)=\log(200K)+\log(\rho/100)\) gives exactly \(J\).
The remaining collision calculation gives the coefficient \(c=\log2/4\)
and remainder \eqref{eq:collisionremainder}. We retain the proposition
as an inherited analytic input; these observations identify the growth
quantities being passed to it.

\subsection{Finite row translation}
For each row \(r=(j,s,\beta)\), write
\[
 Z_i+jr_{\sigma i}+L_i(t)
 =Z_i+j\omega+\log(1+t)
       +j(r_{\sigma i}-\omega)-\tau_i(t).
\]
Expand in the constant errors and tails, and extract \(t^sZ^\beta\).
Only finitely many coefficients can contribute. A choice for that row
has an assigned transverse index \(\alpha_r\), a remaining order
\(\ell_r\le s\), and a scalar \(A_r\). The inherited translation
identities and coefficient bounds give
\begin{equation}\label{eq:rowbound}
 |A_r|\le\exp\left(-\nu(\wt{\alpha_r}-\wt{\beta_r})
                         +H\left(T+\frac{\log2}{w_*}\right)\right).
\end{equation}
A nonzero scalar has \(\alpha_r\ge\beta_r\) componentwise, hence
\(\wt{\alpha_r}\ge\wt{\beta_r}\). The assigned indices can be
restricted to \(\wt\alpha\le H\): larger indices have zero
coefficients in every column by the column degree bound. The tail begins
at \(T_i\), and \(v_0T_i\ge F_0w_i\) is the order budget used to
obtain the \(\nu/F_0\) term in \(T\). The binomial and coefficient
bounds give the other terms displayed in \eqref{eq:errors}.

More explicitly, multilinearity in the rows gives the finite identity
\begin{equation}\label{eq:detexpansion}
 \Delta_\sigma=\sum_f\left(\prod_r A_{r,f}\right)\det B_f,
\end{equation}
where \(f\) chooses one translation term for every row, and \(B_f\)
has the entries \eqref{eq:collisionmatrix}. Its orders satisfy
\(\sum_r\ell_r\le\sum_rs_r\le MH/v_0\), so
Proposition~\ref{prop:collision} applies to each term. The number of
choices per row is bounded by
\begin{equation}\label{eq:choicecount}
 Q_H=(\lfloor H\rfloor+1)^{2m}
                    (\lfloor H/v_0\rfloor+1).
\end{equation}
This bound follows by counting the transverse and auxiliary
multi-indices and the tail order. We use it only as a polynomial bound
for fixed parameters. Equations~\eqref{eq:rowbound}--\eqref{eq:choicecount}
are the inherited finite translation interface, with the added affine
error \(\log2\) from \eqref{eq:affineerror}. No infinite determinant
expansion is being interchanged with a limit.

\subsection{Two ways for a translation term to be small}
Define the low-index count and the collision rate by
\begin{equation}\label{eq:capacity}
 N_H=\#\{\alpha\in\N^m:\wt\alpha\le AH\},\qquad
 C_H=\frac{c\eta^2M}{HN_H}.
\end{equation}
The zero index is included, so \(N_H\ge1\). Fix one nonzero term in
\eqref{eq:detexpansion}, and put \(W=\sum_r\wt{\alpha_r}\).
Multiplying \eqref{eq:rowbound} over the rows gives the scalar estimate
\begin{equation}\label{eq:scalarproduct}
 \left|\prod_r A_{r,f}\right|
 \le\exp\left(-\nu(W-MHb_H)+
                  MH\left(T+\frac{\log2}{w_*}\right)\right).
\end{equation}
We have \(W\ge MHb_H\).

If \(W\ge A(1-\eta)MH\), the first term in the exponent of
\eqref{eq:scalarproduct} is at most
\(-\nu(A(1-\eta)-b_H)MH\). We may discard the nonpositive
collision term in \eqref{eq:collisionbound}.

If \(W<A(1-\eta)MH\), fewer than \((1-\eta)M\) of the assigned
indices can have weight greater than \(AH\). At least \(\eta M\)
assignments therefore fall among the \(N_H\) low indices. If their
multiplicities are \(n_\alpha\), Cauchy--Schwarz gives
\[
 \eta^2M^2\le\left(\sum_{\wt\alpha\le AH}n_\alpha\right)^2
 \le N_H\sum_{\wt\alpha\le AH}n_\alpha^2
 \le N_H\sum_\alpha n_\alpha^2.
\]
The collision saving is consequently at least \(MHC_H\). In this
case we may discard the translation saving, since \(W-MHb_H\ge0\).

Thus every term has one of two savings; we do not claim that every term
has both. Taking the larger of their exponential upper bounds and using
the at most \(Q_H^M\) terms in \eqref{eq:detexpansion} proves
\begin{proposition}[Bound at a good embedding]\label{prop:goodbound}
For every good embedding and every nonzero selected minor,
\begin{equation}\label{eq:goodbound}
 \frac{\log|\Delta_\sigma|}{MH}
 \le V+\frac{\log2}{w_*}+r_a(H)
       +\max\{-C_H,-\nu(A(1-\eta)-b_H)\},
\end{equation}
where
\[
 r_a(H)=r_{\rm coll}(M(H),H)+\frac{\log Q_H}{H}
 \longrightarrow0.
\]
\end{proposition}

To justify this last limit and the collision limit, the fixed-weight
simplex counts give
\begin{align}
 M(H)&\sim\frac{K\theta^m}{(m+1)!v_0\prod_iw_i}\,H^{m+1},
 \label{eq:counts}\\
 N_H&\sim\frac{A^m}{m!\prod_iw_i}\,H^m.\notag
\end{align}
For positive fixed weights these are the usual lattice counts, obtained
by rescaling the simplex and comparing the lattice cubes along its
boundary. Strict versus non-strict inclusion changes only the lower-order
boundary contribution. In particular \(\log M(H)/H\to0\), and
\eqref{eq:choicecount} gives \(\log Q_H/H\to0\). Dividing the two
asymptotics in \eqref{eq:counts} also gives
\begin{equation}\label{eq:ratelimit}
 C_H\longrightarrow
 \frac{c\eta^2K\theta^m}{(m+1)v_0A^m}=C_\infty.
\end{equation}
No uniformity in the dimension or approximation weights is needed; those
choices have already been made.

\section{The remaining embeddings and the multiplier cost}\label{sec:coarse}
At any embedding, including one outside \(G\),
\eqref{eq:approxdata} gives \(|p_i/q_i|\le|x|+1\). Hence
\begin{equation}\label{eq:coarsecenters}
 |r_{\sigma i}|\le|\sigma(u)|(|x|+1+|\sigma(v)|)\le R.
\end{equation}
An entrywise coefficient estimate is enough here. On the circles
\(|t|=1/2\), \(|Z_i|=1\), we have
\[
 |L_i(t)|\le\sum_{k\ge1}\frac{2^{-k}}k\le1,
 \quad |1+t|\le3/2,
 \quad |Z_i+jr_{\sigma i}+L_i(t)|\le KR+2.
\]
Cauchy's coefficient estimate therefore bounds an entry in
\eqref{eq:matrix} by
\[
 2^s\exp(\kappa Kh)(3/2)^h(KR+2)^{\sum_i\gamma_i}.
\]
The row and column budgets give \(s\le H/v_0\), \(h\le H/w_0\),
and \(\sum_i\gamma_i\le H/w_*\). Increasing \(KR+2\) to
\(2(KR+2)\) makes this at most \(\exp(HU(R))\), with \(U\) from
\eqref{eq:errors}. Expanding the determinant into \(M!\) permutation
terms yields the following bound.

\begin{proposition}[Coarse determinant bound]\label{prop:coarse}
At every embedding, for a nonzero selected minor,
\begin{equation}\label{eq:coarse}
 \frac{\log|\Delta_\sigma|}{MH}
 \le U(R)+r_c(H),\qquad r_c(H)=\frac{\log M(H)}H\longrightarrow0.
\end{equation}
\end{proposition}
Indeed \(\log(M!)/(MH)\le\log M/H\), and the limit follows from
\eqref{eq:counts}. The bound is uniform in the chosen columns and in the
embedding; it does not need \(e^\omega=z_\sigma\).

There remains the extra \(\log2/w_*\) in the good estimate. It is
small, but fixed as \(H\) grows, so it cannot be included in \(r_a(H)\).
Instead, for \(K\ge1\) and \(R\ge0\),
\[
 2\cdot2(KR+2)\le2\bigl(K(2R+2)+2\bigr).
\]
Taking logarithms gives
\begin{equation}\label{eq:absorb}
 U(R)+\frac{\log2}{w_*}\le U(2R+2)=U(R').
\end{equation}
This is why the parameter construction used \(R'\) from the outset.
The multiplier has a cost, and this inequality is where we pay it.

\section{The determinant contradiction}\label{sec:contradiction}
We can now complete the proof of Theorem~\ref{thm:common}. For a height
at which interpolation supplies a nonzero selected minor, apply
Proposition~\ref{prop:goodbound} at the two good embeddings and
Proposition~\ref{prop:coarse} at the other two. Lemma~\ref{lem:normcomparison}
and \eqref{eq:absorb} give
\begin{equation}\label{eq:combined}
 -2(1-b_H)-2E-U(R')-r_c(H)
 \le V+r_a(H)+
             \max\{-C_H,-\nu(A(1-\eta)-b_H)\}.
\end{equation}
Put
\begin{equation}\label{eq:SH}
 S_H=2E+U(R')+V+r_a(H)+r_c(H).
\end{equation}
The strict margins \eqref{eq:twomargins}, the limits of the two
remainders, and \eqref{eq:ratelimit} imply that, for all sufficiently
large \(H\),
\begin{equation}\label{eq:eventualmargins}
 S_H<g,\qquad 2+S_H<C_H.
\end{equation}

Suppose first that the maximum in \eqref{eq:combined} is the
approximation term. Rearranging gives
\[
 \nu(A(1-\eta)-b_H)-2(1-b_H)\le S_H<g.
\]
But, using \eqref{eq:gap}, its left side is
\[
 g+(\nu-2)(\theta-b_H)\ge g,
\]
because \(\nu>4\) and \(b_H\le\theta\). This is a contradiction.

Suppose instead that the maximum is the collision term. Then
\eqref{eq:combined} gives
\[
 C_H\le2(1-b_H)+S_H\le2+S_H,
\]
contradicting the second inequality in \eqref{eq:eventualmargins}.
If the two terms are equal, either argument applies.

We have shown that full row rank is impossible at every sufficiently
large height. Lemma~\ref{lem:cofinal} supplies full row rank at a cofinal
set of heights, which is incompatible with that conclusion. Thus
\eqref{eq:bad} is false for the fixed \(\nu>4\). The logical reduction
in Section~\ref{sec:conventions} proves \eqref{eq:eventual} for every
\(\nu>4\), and Lemma~\ref{lem:exponent} proves irrationality and the
exponent bound. This completes the proof of Theorem~\ref{thm:common}.

\section{The three concrete constructions}\label{sec:examples}
Use the field
\begin{equation}\label{eq:field}
 F=\Q(\sqrt3)(\sqrt2)=\Q(\sqrt2,\sqrt3).
\end{equation}
Both generators are algebraic integers. For completeness, its degree is
four: \(\Q(\sqrt3)\) has degree two, and if
\((r+s\sqrt3)^2=2\) with \(r,s\in\Q\), comparison of coefficients
gives \(2rs=0\). The case \(s=0\) would make two a rational square.
The case \(r=0\) gives \(3s^2=2\), hence \((3s)^2=6\), making six
a rational square. Both are impossible. Adjoining \(\sqrt2\) therefore
doubles the degree.

Its four embeddings are the real maps
\begin{equation}\label{eq:signs}
 \sigma_{\varepsilon,\tau}(\sqrt2)=\varepsilon\sqrt2,
 \qquad \sigma_{\varepsilon,\tau}(\sqrt3)=\tau\sqrt3,
 \qquad \varepsilon,\tau\in\{+1,-1\}.
\end{equation}
Each application uses the distinguished embedding \(\sigma_{+,+}\).
The data and good embeddings are summarized below; we verify the
compatibility identities individually.

\begin{center}
\begin{tabular}{@{}c c c c l@{}}
\toprule
Target & \(a\) & \(u\) & \(v\) & Good embeddings\\
\midrule
\(x_1\) & \(2\) & \(1\) & \(\sqrt2\) & \(\sigma_{+,+},\sigma_{+,-}\)\\
\(x_2\) & \(4\) & \(\sqrt2\) & \(0\) & \(\sigma_{+,+},\sigma_{+,-}\)\\
\(x_3\) & \(3+2\sqrt2\) & \(\sqrt2\) & \(\sqrt3\) & \(\sigma_{+,+},\sigma_{-,+}\)\\
\bottomrule
\end{tabular}
\end{center}

\subsection{The algebraic shift}
Take \(a=2\), \(u=1\), \(v=\sqrt2\), and
\(x=x_1=\log2+\sqrt2\). These elements are integral, \(x_1>0\), and
the distinguished value of \(a\) is \(2>0\), different from one.
At either good embedding,
\[
 \sigma(u)(x_1-\sigma(v))=\log2,
 \qquad \exp(\log2)=2=\sigma(a).
\]
Also \(|\sigma(u)|=1\le2\). The two good embeddings are distinct
because they have different values on \(\sqrt3\). All hypotheses of
Theorem~\ref{thm:common} hold, proving \eqref{eq:eventual} for \(x_1\)
and hence Theorem~\ref{thm:shift}.

\subsection{The algebraic multiple}
Take \(a=4\), \(u=\sqrt2\), \(v=0\), and
\(x=x_2=\sqrt2\log2\). Again all the algebraic data are integral,
the target is positive, and the distinguished base is positive and not
one. At either good embedding,
\[
 \sigma(u)(x_2-\sigma(v))=\sqrt2\,\sqrt2\log2=2\log2,
 \qquad \exp(2\log2)=4=\sigma(a).
\]
The multiplier bound is \(|\sigma(u)|=\sqrt2\le2\).
Theorem~\ref{thm:common} gives \eqref{eq:eventual} for \(x_2\), and
Lemma~\ref{lem:exponent} gives Theorem~\ref{thm:scale}.

In these first two constructions the algebraic data lie in the smaller
field \(\Q(\sqrt2)\). The auxiliary \(\sqrt3\) supplies two distinct
embeddings with the required values. Their repeated images do not erase
them from the norm: they are still distinct embeddings of the degree-four
field in which the common theorem is applied.

\subsection{The shifted quadratic period}
Put \(b=3+2\sqrt2>1\), and take \(a=3+2\sqrt2\),
\(u=\sqrt2\), \(v=\sqrt3\). The target is
\(x_3=\log b/\sqrt2+\sqrt3>0\). At \(\sigma_{+,+}\),
\[
 \sigma_{+,+}(u)(x_3-\sigma_{+,+}(v))=\log b,
\]
so the exponential is \(b=\sigma_{+,+}(a)\).
At \(\sigma_{-,+}\), the same computation gives \(-\log b\).
Now
\begin{equation}\label{eq:unit}
 (3+2\sqrt2)(3-2\sqrt2)=1,
\end{equation}
so
\[
 \exp(-\log b)=b^{-1}=3-2\sqrt2=\sigma_{-,+}(a).
\]
Both good embeddings have multiplier absolute value \(\sqrt2\le2\),
and all three field elements are integral. Theorem~\ref{thm:common}
therefore proves \eqref{eq:eventual} for \(x_3\), and proves
Theorem~\ref{thm:period}.

Here the two good embeddings change the sign of the normalized logarithmic
coordinate while keeping the shift \(\sqrt3\) fixed. Identity
\eqref{eq:unit} is what makes their exponential bases match. No
compatibility is needed at the embeddings that change the sign of the
shift; their contribution has already been covered by the coarse bound.

\section{Scope of the conclusions}\label{sec:scope}
For each of our three numbers, the conclusion is
\(2\le\mu(x_j)\le4\). The proof does not assert \(\mu(x_j)=4\),
and it gives no bound with exponent exactly four. It also gives no
numerical value for the threshold in \eqref{eq:eventual}. We make no
claim here about an algorithm for obtaining such a threshold.

The shared theorem has deliberately concrete hypotheses: degree four,
integral \(a,u,v\), a positive distinguished base different from one, and two
compatible embeddings with multipliers bounded by two. A more general
degree-to-embedding ratio may suggest a similar calculation, but it would
need its own parameter and arithmetic argument. Likewise, extending the
statement to nonintegral affine data would introduce denominator costs
that are absent here. Neither extension is asserted by the present paper.

The formalization complements this argument by checking the specified
theorems and the dependencies used to prove them. It does not establish
priority, replace mathematical exposition, or certify statements outside
the checked endpoints. The source and verification record are described
next so that readers can identify the exact scope of those checks.

\appendix
\section{Formalization and verification}\label{app:formalization}
The accompanying package is \path{formalizations/AffineLogarithms} in
the source repository~\cite{affine}. The proof sources are frozen; this
manuscript explains those sources rather than changing their statements.
The common endpoint is
\begin{center}\small
\path{OAI.PairedAffine.Endpoint.paired_eventualLowerBound}.
\end{center}
Its parameters are exactly the field elements, integrality proofs,
degree, distinguished embedding, real target, and good-embedding
conditions in Theorem~\ref{thm:common}. The implementation uses
\(\exp x\) as a positive real proxy in an inherited data structure, so
that its logarithm is \(x\). This proxy is not the algebraic base
\(a\), and no algebraicity is required of it.

\subsection{Statement correspondence}
The twelve public declarations consist of four statements for each
target: the packaged eventual bound, its expanded quantified form,
irrationality, and the exponent upper bound. Three further comparison
wrappers state irrationality and exclude Mathlib's \texttt{LiouvilleWith}
\(\nu\,x\) for every \(\nu>4\), giving fifteen literal comparison targets.
With the common namespace
\texttt{OAI.PairedAffine}, the target stems are
\begin{center}\small
\begin{tabular}{@{}ll@{}}
\toprule
Theorem & Declaration stem\\
\midrule
\ref{thm:shift} & \texttt{log\_two\_add\_sqrt\_two}\\
\ref{thm:scale} & \texttt{sqrt\_two\_mul\_log\_two}\\
\ref{thm:period} & \texttt{shifted\_period}\\
\bottomrule
\end{tabular}
\end{center}
The suffixes are \texttt{\_eventualLowerBound},
\texttt{\_explicit\_bound}, \texttt{\_irrational}, and
\path{_irrationalityExponent_le_four}. In the second suffix,
``explicit'' refers to an expanded quantified statement, not to an
explicit numerical threshold. The file \texttt{evidence/target-map.json}
records the literal expressions and comparison targets.

The exponent convention is in
\path{OAI/NumberTheory/PiExponent/Statement.lean}; its use in this
package is justified by \path{AlgebraicLog/ExponentConsequence.lean}.
The standard comparison wrappers are in \path{Solution.lean};
\path{StandardBridge.lean} supplies the implication from the eventual
bound to the exclusion of \texttt{LiouvilleWith}.
In particular, the comparison is not merely between two names with a
similar intended meaning.

\subsection{Where the main steps are checked}
The following source map gives the principal interfaces. Paths are
relative to the package root; declaration names below omit the initial
\texttt{OAI.}.

\begin{longtable}{@{}>{\raggedright\arraybackslash}p{0.27\textwidth}>{\raggedright\arraybackslash}p{0.67\textwidth}@{}}
\toprule
Mathematical step & Source and declaration\\
\midrule
\endhead
Shape and parameter margins &
\path{AlgebraicLog/ParameterShape.lean},
\texttt{exists\_quadratic\_shape};
\path{AlgebraicLog/DimensionBudget.lean};
\path{AlgebraicLog/AdmissibleParameters.lean},
\path{exists_admissible_parameters},
\texttt{error\_sum\_lt\_gap}, and
\path{collision_exceeds_error_sum}.\\[4pt]
Weighted interpolation &
\path{Logarithm/Interpolation.lean},
\path{eventually_weighted_logarithmic_interpolation};
its hypotheses are the fields of
\path{Logarithm/GeometryData.lean}.\\[4pt]
Truncation and full row rank &
\path{Logarithm/ScaledTruncation.lean},
\path{formalLog_packets_surjective_iff_truncated};
\path{Logarithm/ScaledMatrix.lean},
\path{truncatedLogMatrix_surjective_of_formalLog_packets};
\path{PairedAffine/Geometry.lean},
\texttt{cofinal\_matrix}.\\[4pt]
Integral entries and selected minor &
\path{PairedAffine/ArithmeticClearing.lean},
\path{exists_integral_entry_lift};
\path{PairedAffine/Arithmetic.lean},
\path{exists_integral_selected_minor} and
\texttt{log\_selectedScalar\_le}.\\[4pt]
Norm comparison &
\path{PairedAffine/NormSelected.lean},
\path{paired_selected_normalized_bound}.\\[4pt]
Collision estimate &
\path{Logarithm/ScaledCollision.lean},
\path{formal_scaled_collision_bound}; the analytic foundation
is in \path{OAI/NumberTheory/PiExponent/Analysis/}.\\[4pt]
Finite translation and good bound &
\path{PairedAffine/Analytic.lean},
\texttt{norm\_choiceScalar\_le}, \texttt{minor\_expansion},
and \texttt{minor\_analytic\_bound}.\\[4pt]
Coarse bound and multiplier &
\path{PairedAffine/AnalyticCoarse.lean},
\texttt{minor\_log\_bound} and
\path{conjugateCost_add_log_two_le}.\\[4pt]
Contradiction and endpoint &
\path{AlgebraicLog/ContradictionNumerics.lean},
\path{quadratic_determinant_bounds_inconsistent};
\path{PairedAffine/Endpoint.lean},
\texttt{no\_fixed\_data} and \texttt{paired\_eventualLowerBound}.\\[4pt]
Concrete field and targets &
\path{PairedAffine/Field.lean},
\path{PairedAffine/AffineExamples.lean},
\path{PairedAffine/Concrete.lean}, and
\path{PairedAffine/Main.lean}.\\
\bottomrule
\end{longtable}

\subsection{Recorded checks and their limits}
The primary toolchain is Lean~4.34.1, with Mathlib pinned to
\path{d13f23b723b8a846827a245b89c10fc7d3f11612}. The package records
source compilation, axiom inspection, literal statement comparison,
export identity, and proof-body replay as distinct checks. The permitted
foundational axioms are \texttt{propext}, \texttt{Classical.choice},
and \texttt{Quot.sound}. The solution does not import the challenge
file's statement placeholders.

The retained original export passed the fifteen statement comparisons,
the Lean~4.34.1 checker, and Rust Nanoda. Nanoda supplies an independently
implemented kernel. A supplemental Lean~4.34.0 reader also replayed that
export; it belongs to the same kernel family as Lean and lacks the later
runtime fixes in 4.34.1, so it is historical supplementary evidence rather
than the primary Lean check.

A fresh hosted run on 11 October 2026, run
\href{https://github.com/Mattie/math/actions/runs/38139172013}{38139172013},
attempt~1, used repository commit
\path{b15a6431eca8632a48c1cd84e5200edd30f0670c}. The retained audit
records all fifteen comparison targets as passed, with Lean and Rust
Nanoda, a fresh build of the requested candidate dependency closure, and
checks of the 969 frozen source/specification inputs. Official Mathlib
dependency caches remained a declared trust input. The run did not retain
a new standalone export; the earlier export and the fresh hosted results
are separate evidence records.

The checks include acceptance and rejection controls, including statement
mismatch and imported-admission cases. These are bounded tests, not a
proof that the tools have no bugs or that process isolation contains
every hostile program. The package's \path{CHECKERS.md} discusses the
selected checker revisions, including structured axiom-name checks
needed around known Nanoda limitations. Its \path{VERIFY.md},
\path{PROVENANCE.md}, and
\path{evidence/hosted-result-2026-10-11.json} distinguish the original
replay, package checks, and fresh hosted run. The verification claim here
is acceptance of specified statements under that recorded toolchain and
axiom policy, not human peer review.

\section{Source attribution and preparation}\label{app:attribution}
The inherited source pin is
\path{0cfe10002ea95c45721b93d18bffbe2c3bbbcbbd} in the
\texttt{Mattie/math} repository. It incorporates OpenAI's
\texttt{OAI.NumberTheory.PiExponent} library at revision
\path{adc7f1241b42e322a6451854ab7e4b4c146bf78a}. The rational-logarithm
and quadratic-logarithm extensions supply mathematical constructions and
estimates used throughout this paper. The original affine work extends
that framework to the common integral affine data and the paired
embedding argument. We have stated the inherited interfaces explicitly
to avoid presenting those foundations as if they were established here.

AI assistance, including OpenAI Codex agents, was used in preparing the
manuscript and accompanying formalization. The author is responsible for
the submitted claims and their presentation. The formal verification
results describe the checked proof artifacts; AI assistance itself is
not evidence for a mathematical claim.

The package retains the inherited OpenAI library's Apache~2.0 license and
notices. Original \texttt{PairedAffine}, \texttt{AlgebraicLog}, and
\texttt{Logarithm} sources and associated original code are distributed
under the MIT license; original explanatory prose is under CC~BY~4.0.
Lean, Mathlib, and the other inherited tools retain their respective
notices. The upstream authors and organizations are contributors to the
dependencies, not endorsers of this paper.

\begin{thebibliography}{9}

\bibitem{luchin}
M.~Yu.~Luchin and V.~Kh.~Salikhov,
\emph{Approximating \(\ln2\) by numbers in the field \(\Q(\sqrt2)\)},
Izvestiya: Mathematics \textbf{82} (2018), no.~3, 549--577.
Theorem~1, p.~549.
\url{https://www.mathnet.ru/php/getFT.phtml?jrnid=im&option_lang=eng&paperid=8518&what=fullteng}.

\bibitem{openai}
OpenAI,
\emph{\texttt{OAI.NumberTheory.PiExponent}: formal source library},
revision \path{adc7f1241b42e322a6451854ab7e4b4c146bf78a}.
\url{https://github.com/openai/math/tree/adc7f1241b42e322a6451854ab7e4b4c146bf78a}.

\bibitem{rational}
R.~M.~Casper,
\emph{The irrationality exponent of positive rational logarithms is 2},
preprint and Lean source package, 7 October 2026.
Inherited source snapshot at repository revision
\path{0cfe10002ea95c45721b93d18bffbe2c3bbbcbbd}, directory
\path{preprints/The-irrationality-exponent-of-positive-rational-logarithms-is-2-October-7-2026/lean}.
\url{https://github.com/Mattie/math/tree/0cfe10002ea95c45721b93d18bffbe2c3bbbcbbd}.

\bibitem{quadratic}
R.~M.~Casper,
\emph{The irrationality exponent of positive quadratic logarithms is at most 4},
preprint and Lean source package, 8 October 2026.
The frozen \texttt{AlgebraicLog} modules used here are included in the
affine package~\cite{affine} and identified by its source manifest.

\bibitem{affine}
R.~M.~Casper,
\emph{AffineLogarithms: paired affine logarithms, proof sources and
verification records}, October 2026.
Directory \path{formalizations/AffineLogarithms} in
\url{https://github.com/Mattie/math/tree/b15a6431eca8632a48c1cd84e5200edd30f0670c}.
Frozen mathematical sources, statement map, and verification provenance;
hosted verification run of 11 October 2026 cited in
Appendix~\ref{app:formalization}.

\end{thebibliography}
\end{document}
